% --- 0걸음 → 1걸음
\begin{frame}{0걸음 $\to$ 1걸음: 첫 선은 어떻게 생기나}
\vfill
\begin{columns}[c]
\begin{column}{0.40\textwidth}\centering\includegraphics[width=\linewidth]{sv_step1.png}\end{column}
\begin{column}{0.59\textwidth}
{\large \textbf{0걸음}: 값 모두 0 $\to$ \kq{20개 모두} 가장 가깝다}\\[0.5em]
{\large \textbf{1걸음}: 빨강은 \kq{더하고}, 파랑은 \kb{뺀다}}\\[0.3em]
\renewcommand{\arraystretch}{1.1}
{\normalsize\begin{tabular}{l|rr}
 & $x_1$ 합 & $x_2$ 합\\ \hline
빨강 10개 & $8.3$ & $9.1$\\
파랑 10개 & $-8.2$ & $-9.1$\\ \hline
빨강 $-$ 파랑 & $16.5$ & $18.2$
\end{tabular}}\\[0.4em]
{\normalsize $\mathbf{w} = \eta \sum_{i=1}^{20} y_i\mathbf{x}_i = \eta\big(\sum_{\text{빨강}}\mathbf{x}_i - \sum_{\text{파랑}}\mathbf{x}_i\big)$}\\[0.1em]
{\normalsize $b = \eta \sum_{i=1}^{20} y_i = \eta\,(10 - 10) = 0$}\\[0.3em]
{\large $\mathbf{w} = 0.05 \times (16.5,\ 18.2) = (0.83,\ 0.91)$}\\[0.2em]
{\large $\lVert\mathbf{w}\rVert = 1.23$ 로 normalize 하면 \kb{$(0.67,\ 0.74)$}}
\end{column}
\end{columns}
\vfill
\note{0걸음부터 따라가 봅시다. 처음 w는 0, 0, b는 0이라 선이 없습니다. 어느 점을 넣어도 값이 0이라 20개가 모두 똑같이 가장 가까운 점입니다. 그래서 20개가 모두 밉니다. 빨간 점은 자기 좌표를 더하고 파란 점은 뺍니다. 빨강 10개의 좌표 합은 8.3, 9.1, 파랑은 마이너스 8.2, 마이너스 9.1, 빼면 16.5, 18.2입니다. 학습률 0.05를 곱하면 0.83, 0.91, 길이가 1.23으로 1보다 크니 1.23으로 나눠 0.67, 0.74가 됩니다. b는 빨강 10개 빼기 파랑 10개, 0입니다. 앞에서 본 가장 쉬운 방법, 20개의 y i x i 평균과 같은 방향이고 b도 0이라, 첫 선은 그 평균 선과 똑같습니다. 0에서 출발해서 20개가 모두 같은 값일 때, 값이 같은 점은 다 같이 민다는 이 덱의 규칙을 쓰면 이렇게 됩니다. 이건 교과서에 있는 정리가 아니라 이 방법으로 직접 계산한 결과입니다. m은 0.62로, 이미 답 0.67에 가깝습니다.}
\end{frame}

